% GENERATED FILE. Edit canonical lessons, then run npm run generate:books.
% canonical-source-hash: fnv1a64:11ea2186effcf95a
% canonical-lessons: UR-C278-darogha, UR-C278-riport, UR-C278-darj-karna, UR-C278-havalat, UR-C278-bayan

\chapter{At the Police Station}
\label{ch:ur-at-the-police-station}
\hlchaptermodality{278}

\begin{chapteropening}
\textbf{By the end of this chapter:} \emph{I can report a crime at the police station.}

Complete the last lesson of chapter 278: I can report a crime at the police station.
\end{chapteropening}

\section[\texorpdfstring{\ur{داروغہ} (dāroghā)}{داروغہ (dāroghā)}]{\ur{داروغہ} (dāroghā) — a police inspector}

\label{lesson:UR-C278-darogha}



\begin{quote}
Before the new one: say the Urdu for on the run, then the Urdu for a criminal.
\end{quote}

\subsection*{You'll want to know: \ur{داروغہ}}
\textbf{\ur{داروغہ}} — \emph{dāroghā} — \textquotedblleft{}a police inspector\textquotedblright{}.

\begin{quote}
\ur{داروغہ} \ur{تھانے} \ur{میں} \ur{نہیں} \ur{تھا۔}
\end{quote}

\begin{quote}
\emph{dāroghā thāne meṅ nahīṅ thā.}
\end{quote}

\begin{quote}
The inspector was not at the police station.
\end{quote}

Masculine.

\begin{cousinweb}
is usually traced to Persian \emph{dārogha}, \textquotedblleft{}an overseer\textquotedblright{}, a word Persian is thought to have taken from Mongolian. In a town the \textbf{\ur{داروغہ}} runs the \textbf{\ur{تھانہ}}.
\end{cousinweb}

\begin{grammarlens}[title={What you've built}]
One more word for telling what happened.
\end{grammarlens}

\subsection*{Your turn}
\begin{itemize}
\raggedright
  \item \emph{Say it:} \emph{dāroghā}
  \item \emph{Say it:} \emph{dāroghā}, once more
  \item \emph{Say it:} \emph{dāroghā thāne meṅ nahīṅ thā.} — The inspector was not at the police station.
  \item \emph{Recall:} say the Urdu for on the run, then the Urdu for a criminal, then say \emph{dāroghā} again
\end{itemize}

\begin{tcolorbox}[breakable,colback=teal!4,colframe=teal!35!black,title={Before you move on}]
What does \ur{داروغہ} mean? (\textquotedblleft{}A police inspector\textquotedblright{}.) Say it once more.
\end{tcolorbox}
\section[\texorpdfstring{\ur{رپورٹ} (riporṭ)}{رپورٹ (riporṭ)}]{\ur{رپورٹ} (riporṭ) — a report (to the police)}

\label{lesson:UR-C278-riport}



\begin{quote}
Before the new one: say the Urdu for a criminal, then the Urdu for a police inspector.
\end{quote}

\subsection*{You'll want to know: \ur{رپورٹ}}
\textbf{\ur{رپورٹ}} — \emph{riporṭ} — \textquotedblleft{}a report (to the police)\textquotedblright{}.

\begin{quote}
\ur{پولیس} \ur{نے} \ur{رپورٹ} \ur{لکھی۔}
\end{quote}

\begin{quote}
\emph{polīs ne riporṭ likhī.}
\end{quote}

\begin{quote}
The police wrote up a report.
\end{quote}

Feminine.

\begin{cousinweb}
\textbf{\ur{رپورٹ}} (\emph{riporṭ}) is borrowed from English \emph{report}. At the police station it is the written account of a crime, and it is feminine: \textbf{\ur{رپورٹ} \ur{لکھی}}.
\end{cousinweb}

\begin{grammarlens}[title={What you've built}]
One more word for telling what happened.
\end{grammarlens}

\subsection*{Your turn}
\begin{itemize}
\raggedright
  \item \emph{Say it:} \emph{riporṭ}
  \item \emph{Say it:} \emph{riporṭ}, once more
  \item \emph{Say it:} \emph{polīs ne riporṭ likhī.} — The police wrote up a report.
  \item \emph{Recall:} say the Urdu for a criminal, then the Urdu for a police inspector, then say \emph{riporṭ} again
\end{itemize}

\begin{tcolorbox}[breakable,colback=teal!4,colframe=teal!35!black,title={Before you move on}]
What does \ur{رپورٹ} mean? (\textquotedblleft{}A report (to the police)\textquotedblright{}.) Say it once more.
\end{tcolorbox}
\section[\texorpdfstring{\ur{درج} \ur{کرنا} (darj karnā)}{درج کرنا (darj karnā)}]{\ur{درج} \ur{کرنا} (darj karnā) — to register (in a record)}

\label{lesson:UR-C278-darj-karna}



\begin{quote}
Before the new one: say the Urdu for a police inspector, then the Urdu for a report (to the police).
\end{quote}

\subsection*{You'll want to know: \ur{درج} \ur{کرنا}}
\textbf{\ur{درج} \ur{کرنا}} — \emph{darj karnā} — \textquotedblleft{}to register (in a record)\textquotedblright{}.

\begin{quote}
\ur{داروغے} \ur{نے} \ur{رپورٹ} \ur{درج} \ur{کی۔}
\end{quote}

\begin{quote}
\emph{dāroghe ne riporṭ darj kī.}
\end{quote}

\begin{quote}
The inspector registered the report.
\end{quote}

Past, with \ur{نے} (he or she …): \ur{اس} \ur{نے} \ur{رپورٹ} \ur{درج} \ur{کی}, \emph{us ne riporṭ darj kī}.

\begin{cousinweb}
\textbf{\ur{درج} \ur{کرنا}} (\emph{darj karnā}) is built from \emph{darj}, usually traced, through Persian, to Arabic \emph{darj}, \textquotedblleft{}a roll, a written record\textquotedblright{}, and \textbf{\ur{کرنا}}: to put a thing on the record. The thing written agrees: \textbf{\ur{رپورٹ} \ur{درج} \ur{کی}}.
\end{cousinweb}

\begin{grammarlens}[title={What you've built}]
One more word for telling what happened.
\end{grammarlens}

\subsection*{Your turn}
\begin{itemize}
\raggedright
  \item \emph{Say it:} \emph{darj karnā}
  \item \emph{Say it:} \emph{darj karnā}, once more
  \item \emph{Say it:} \emph{dāroghe ne riporṭ darj kī.} — The inspector registered the report.
  \item \emph{Recall:} say the Urdu for a police inspector, then the Urdu for a report (to the police), then say \emph{darj karnā} again
\end{itemize}

\begin{tcolorbox}[breakable,colback=teal!4,colframe=teal!35!black,title={Before you move on}]
What does \ur{درج} \ur{کرنا} mean? (\textquotedblleft{}To register (in a record)\textquotedblright{}.) Say it once more.
\end{tcolorbox}
\section[\texorpdfstring{\ur{حوالات} (havālāt)}{حوالات (havālāt)}]{\ur{حوالات} (havālāt) — a lock-up}

\label{lesson:UR-C278-havalat}



\begin{quote}
Before the new one: say the Urdu for a report (to the police), then the Urdu for to register.
\end{quote}

\subsection*{You'll want to know: \ur{حوالات}}
\textbf{\ur{حوالات}} — \emph{havālāt} — \textquotedblleft{}a lock-up\textquotedblright{}.

\begin{quote}
\ur{چور} \ur{رات} \ur{بھر} \ur{حوالات} \ur{میں} \ur{رہا۔}
\end{quote}

\begin{quote}
\emph{chor rāt bhar havālāt meṅ rahā.}
\end{quote}

\begin{quote}
The thief spent the whole night in the lock-up.
\end{quote}

Feminine.

\begin{cousinweb}
\textbf{\ur{حوالات}} (\emph{havālāt}) is usually traced, through Persian, to Arabic \emph{ḥawāla}, \textquotedblleft{}a handing over, custody\textquotedblright{}, in a plural form: the room at the police station where a person handed over is held for a night or two.
\end{cousinweb}

\begin{grammarlens}[title={What you've built}]
One more word for telling what happened.
\end{grammarlens}

\subsection*{Your turn}
\begin{itemize}
\raggedright
  \item \emph{Say it:} \emph{havālāt}
  \item \emph{Say it:} \emph{havālāt}, once more
  \item \emph{Say it:} \emph{chor rāt bhar havālāt meṅ rahā.} — The thief spent the whole night in the lock-up.
  \item \emph{Recall:} say the Urdu for a report (to the police), then the Urdu for to register, then say \emph{havālāt} again
\end{itemize}

\begin{tcolorbox}[breakable,colback=teal!4,colframe=teal!35!black,title={Before you move on}]
What does \ur{حوالات} mean? (\textquotedblleft{}A lock-up\textquotedblright{}.) Say it once more.
\end{tcolorbox}
\section[\texorpdfstring{\ur{بیان} (bayān)}{بیان (bayān)}]{\ur{بیان} (bayān) — a statement}

\label{lesson:UR-C278-bayan}



\begin{quote}
Before the new one: say the Urdu for to register, then the Urdu for a lock-up.
\end{quote}

\subsection*{You'll want to know: \ur{بیان}}
\textbf{\ur{بیان}} — \emph{bayān} — \textquotedblleft{}a statement\textquotedblright{}.

\begin{quote}
\ur{یہ} \ur{گواہ} \ur{کے} \ur{بیان} \ur{پر} \ur{منحصر} \ur{ہے۔}
\end{quote}

\begin{quote}
\emph{yih gavāh ke bayān par munhasir hai.}
\end{quote}

\begin{quote}
It depends on the witness's statement.
\end{quote}

Masculine.

\begin{cousinweb}
is usually traced, through Persian, to Arabic \emph{bayān}, \textquotedblleft{}a making clear, an account\textquotedblright{}. A \textbf{\ur{گواہ}} gives a \textbf{\ur{بیان}} to the police.
\end{cousinweb}

\begin{grammarlens}[title={What you've built}]
One more word for telling what happened.
\end{grammarlens}

\subsection*{Your turn}
\begin{itemize}
\raggedright
  \item \emph{Say it:} \emph{bayān}
  \item \emph{Say it:} \emph{bayān}, once more
  \item \emph{Say it:} \emph{yih gavāh ke bayān par munhasir hai.} — It depends on the witness's statement.
  \item \emph{Recall:} say the Urdu for to register, then the Urdu for a lock-up, then say \emph{bayān} again
\end{itemize}

\begin{tcolorbox}[breakable,colback=teal!4,colframe=teal!35!black,title={Before you move on}]
What does \ur{بیان} mean? (\textquotedblleft{}A statement\textquotedblright{}.) Say it once more.
\end{tcolorbox}
